% §5.1 From Atom to Set
%
% Cross-reference resolution status:
%
% DEFINED here:
%   \label{sec:from-atom-to-set}                -- this section
%
% RESOLVED (already in source tree from Plans 1-4):
%   \Cref{ch:noisy-bit}                         -- chapter 1
%   \Cref{sec:three-faces}                      -- chapter 1 §1.1
%   \Cref{sec:why-atom}                         -- chapter 1 §1.5
%   \Cref{ch:axioms-kronecker}                  -- chapter 4
%   \Cref{sec:bernoulli-bool-formal}            -- chapter 4 §4.5
%   \Cref{def:bernoulli-bool}                   -- chapter 4 §4.5
%   \Cref{sec:axiom-2}                          -- chapter 4 §4.3
%   \Cref{ax:block-rates-indep}                 -- chapter 4 §4.3
%   \Cref{sec:kronecker}                        -- chapter 4 §4.4
%   \Cref{thm:kronecker}                        -- chapter 4 §4.4
%
% FORWARD REFS within chapter 5 (resolve in Plan 5 Tasks 3-4):
%   \Cref{sec:bloom-from-scratch}               -- §5.2 (Task 3)
%   \Cref{sec:predicate-algebra}                -- §5.3 (Task 4)

\section{From Atom to Set}
\label{sec:from-atom-to-set}

Part~I closed with a promise. \Cref{sec:why-atom} described the noisy
bit as the framework's atom and previewed what Part~II would do with it:
take many of them, arrange them into a structure, and study the algebra
of that structure. The single-Boolean object of chapters 1 through 4
would become a building block, and the natural building was a noisy bit
per universe element. That object is the approximate set, and Part~II
opens with it.

Why is the construction inevitable? A classical set $A \subseteq U$ is
characterized entirely by its membership predicate: a function
$\SetIndicator{A} : \Universe \to \{\top, \bot\}$ that returns $\top$ on
the elements of $A$ and $\bot$ everywhere else. Two sets are equal if
and only if their predicates agree on every input. The predicate is the
set; the set-builder notation $\{x : \SetIndicator{A}(x) = \top\}$ is
just the predicate read backwards. Now replace the exact predicate with
a noisy one. Each query $x \mapsto \tilde A(x)$ returns a Boolean that
is usually but not always the correct membership verdict. Each
evaluation is a noisy bit, an instance of the formal type
$\Bernoulli{\BoolType}$ that \Cref{sec:bernoulli-bool-formal} defined.
An approximate set is the object whose membership predicate has been
made noisy in this way, one query at a time.

\subsection*{The predicate framing}

Let $\Universe$ be a fixed universe (finite or countably infinite) and
let $A \subseteq \Universe$ be the underlying \emph{objective set} we
intend to represent. An \emph{approximate set} of $A$ is a function
\[
  \ASet{A} : \Universe \longrightarrow \Bernoulli{\BoolType}
\]
such that each evaluation $\ASet{A}(x)$ is a $\Bernoulli{\BoolType}$
instance. Per \Cref{def:bernoulli-bool}, a $\Bernoulli{\BoolType}$
instance is a triple $(X, Y, Q)$ that produces a Boolean output when
sampled; we write $\ASet{A}(x) \in \{\top, \bot\}$ when referring to
the sampled output of $\ASet{A}$ at $x$.
The notation is the bookkeeping: $A$ is the truth, $\ASet{A}$
is its noisy representation, and $\ASet{A}(x) = \top$ is read as ``the
structure reports that $x$ is a member of $A$.'' The structure may be a
Bloom filter, a perfect hash filter, a learned classifier with a
threshold, or any other implementation whose membership queries return
Booleans. What makes it an approximate set, in the sense Part~II
develops, is that each per-element query is governed by a single
$Q$-matrix from chapter 4: two rates of error, one rate for the
positives and one rate for the negatives, and nothing else.

Two paragraphs are worth pausing on. First, the predicate framing is
not a definitional choice; it is a recognition. The classical
correspondence between a set and its indicator function is exact, and
the approximate analogue is forced once we ask what a set looks like
when each query may be wrong. The predicate is where the noise lives
and there is nowhere else for it to live, because there is no other
behavioral handle on a set besides the membership test. Second, the
distinction between $A$ and $\ASet{A}$ is the distinction between truth
and report. We never lose access to $A$ at the level of the model: $A$
is the objective set, and the rates of $\ASet{A}$ are defined against
it. What we lose, when we query the structure, is the ability to read
$A$ off directly. We see only the report, $\ASet{A}(x)$, and the rates
tell us how that report relates to the underlying truth.

\subsection*{The two error rates}

Every query $\ASet{A}(x)$ is one of two kinds, determined by whether
$x$ is in $A$. If $x \notin A$, then the correct verdict is $\bot$ and
the structure errs only by reporting $\top$. The false-positive rate
of the approximate set is the conditional probability of that error:
\[
  \fprate \;=\; \Prob{\ASet{A}(x) = \top \Given x \notin A}.
\]
If $x \in A$, then the correct verdict is $\top$ and the structure
errs only by reporting $\bot$. The false-negative rate is the
conditional probability of that error:
\[
  \fnrate \;=\; \Prob{\ASet{A}(x) = \bot \Given x \in A}.
\]
These are the two rates that chapter 1 introduced for a single noisy
bit (\Cref{sec:three-faces}) and that chapter 4 collected into the
channel matrix $Q$ (\Cref{def:bernoulli-bool}). The change in Part~II
is one of grouping rather than of substance. A noisy bit had one
positive input and one negative input; the approximate set has $|A|$
positive inputs (the members) and $|\Universe \setminus A|$ negative
inputs (the non-members). The two-rate description does not grow: the
same $\fprate$ governs every query on a non-member, and the same
$\fnrate$ governs every query on a member.

That collapse from many per-element rates to two per-set rates is not
automatic. In principle every element $x \in \Universe$ could carry
its own rate $\fprate(x)$ or $\fnrate(x)$, and the predicate would
need a function-valued description rather than a two-number one. What
collapses the description is Axiom 2 of chapter 4
(\Cref{ax:block-rates-indep}, \Cref{sec:axiom-2}): the per-block rates
are stable across the block. Membership splits the universe into two
blocks, and Axiom 2 forces a single rate per block. \Cref{sec:predicate-algebra}
makes this precise. For now, the operational takeaway is that an
approximate set is described, at the rate level, by the pair
$(\fprate, \fnrate)$ alone, no matter how large $|\Universe|$ or
$|A|$ is.

\subsection*{The block structure}

The universe $\Universe$ does not enter the rate description as a flat
list of elements. It enters as a two-block partition: the in-set
block $B_1 = A$ and the out-of-set block $B_2 = \Universe \setminus
A$. Membership induces the partition; the partition determines which
rate applies to which query. This is exactly the model-order-two block
structure of chapter 4: two blocks, two rate-rows in the channel
matrix $Q$, one row per block. An approximate set is a
$\Bernoulli{\BoolType}$ instance whose input domain is $\Universe$
rather than $\{\top, \bot\}$, with the membership partition standing
in for the Boolean partition of the atomic case. The matrix $Q$ that
governs each query is the same:
\[
  Q \;=\;
  \begin{pmatrix}
    1 - \fnrate & \fnrate \\
    \fprate & 1 - \fprate
  \end{pmatrix},
\]
with row 1 for the in-set block and row 2 for the out-of-set block.
\Cref{sec:predicate-algebra} writes this out formally and treats the
case where the rates depend on $x$ before Axiom 2 collapses them.

The payoff of the block structure is that the per-query description
generalizes cleanly to many queries. Chapter 4
(\Cref{sec:kronecker}, \Cref{thm:kronecker}) showed that $n$
independent queries on a $\Bernoulli{\BoolType}$ instance produce a
joint distribution that factors as the Kronecker product of the $n$
per-query channels. The same theorem applies here, with the
membership partition in place of the Boolean partition. The Kronecker
factorization of $n$-query behavior also relies on Axiom~1
(\Cref{ax:element-indep}), which we silently invoked above: error
events at distinct elements of $\Universe$ are mutually independent.
Axiom~1 lets us treat the queries as independent draws; Axiom~2 lets
us assign the same per-block rate to every draw within a block.
Together they give the Kronecker product structure of
\Cref{thm:kronecker}. A workload
of $n$ queries on $\ASet{A}$, drawn from the in-set and out-of-set
blocks in some specified ratio, has a joint behavior that factors
exactly as the tensor product of the per-query matrices. The price
of the full joint description is the two rates, not $2^n$
probabilities. The chapter 4 machinery does not need to be rebuilt
for Part~II; it transfers over wholesale, and \Cref{sec:predicate-algebra}
states the transfer precisely.

\subsection*{Where the chapter goes from here}

The construction is so far an abstraction. We have said what an
approximate set is, namely a function from a universe to the noisy-bit
atom of Part~I, and we have named the rates that describe it. We
have not yet said how anyone would build one. The most widely deployed
construction is the Bloom filter, and \Cref{sec:bloom-from-scratch}
develops it from scratch as the canonical instance of the predicate
framing. The construction will have a structural false-negative rate
of zero, a tunable false-positive rate, and a closed-form formula for
the false-positive rate as a function of three integers (the bit-array
size $m$, the hash count $k$, and the number of stored items $n$).
The Bloom filter is the example chapters 6 through 8 keep returning
to, so getting it on the table first matters.

\Cref{sec:predicate-algebra} then develops the predicate algebra in
general. The Bloom filter is one approximate set; the predicate
framing covers every approximate set, and the chapter 4 axioms and
Kronecker theorem apply uniformly across the family.
\Cref{sec:one-sided-two-sided} draws the line between approximate sets
with one structural zero rate and those with two positive rates: a
distinction that organizes chapters 6 and 15.
